\documentclass[a4paper,12pt]{article}

\usepackage{alltt, fancyvrb, url}
\usepackage{graphicx}
\usepackage[utf8]{inputenc}
\usepackage{float}
\usepackage{hyperref}
\usepackage[italian]{babel}
\usepackage[table]{xcolor}
\usepackage{array}
\usepackage{multirow}
\usepackage{titlesec}
\usepackage{listings}
\usepackage[italian]{cleveref}
\usepackage{acronym}

\renewcommand{\thesection}{\arabic{section}}

\sloppy
\setlength{\parindent}{0pt}

\title{\textbf{Project Definition Statement (PDS)}}
\author{}
\date{}

\begin{document}
\maketitle
\tableofcontents

\newpage

\section{Opportunità}
L'azienda ChargeAndTrack gestisce una distribuzione di colonnine di ricarica per auto elettriche in tutta Italia.
E' stata effettuata un'analisi di mercato dei principali distributori di colonnine di ricarica sul territorio
italiano, e si è riscontrato che nessuno di essi offre soluzioni digitali per la ricerca e l'utilizzo di
colonnine di ricarica. Inoltre, è stata condotta un'indagine tra i clienti e si è riscontrato che il 70\% di essi
sarebbe interessato ad una soluzione digitale. In particolare, ha trovato maggiore interesse la possibilità di
pianificare un itinerario di viaggio con soste di ricarica, e la possibilità di ricevere aggiornamenti in tempo reale
sulla ricarica di una propria auto.

\section{Goal}
Si vuole sviluppare un'applicazione software con le caratteristiche emerse dall'indagine con i clienti, in modo da
aumentare il fatturato dell'azienda del 10\% entro un anno dal rilascio dell'applicazione.

\section{Obiettivi}
\begin{enumerate}
    \item Fornire agli utenti un'esperienza personalizzata in base al proprio profilo;
    \item mostrare agli utenti le colonnine di ricarica presenti entro un raggio variabile rispetto alla posizione
    dell'utente e la loro disponibilità in tempo reale;
    \item dare agli utenti la possibilità di avviare/interrompere una ricarica presso una colonnina di una propria auto
    registrata;
    \item permettere agli utenti di ricevere aggiornamenti in tempo reale sulla ricarica di una propria auto;
    \item permettere agli utenti di pianificare un itinerario di viaggio per ottenere un percorso comprendente le soste
    di ricarica necessarie.
\end{enumerate}

\subsection{Rischi, ostacoli ed assunzioni}

Relativamente alla numerazione degli obiettivi, si possono individuare i seguenti rischi, ostacoli ed assunzioni:
\begin{enumerate}
    \item Ostacolo: bisogna mettere in atto misure per garantire la privacy degli utenti;
    \item Rischio: in alcune zone con basso segnale, l'aggiornamento della disponibilità delle colonnine di ricarica
    potrebbe subire dei ritardi;
    \item Assunzione: le colonnine di ricarica sono già dotate dell'hardware necessario per supportare le funzionalità
    richieste dall'applicazione;
\end{enumerate}

\section{Conditions of Satisfaction (CoS)}
\begin{enumerate}
    \item L'applicazione deve essere implementata come web app e app mobile;
    \item l'applicazione deve essere completata entro un anno;
    \item non superare il budget di 200.000 euro;
    \item prestare particolare attenzione alla User Experience (UX);
    \item il percorso generato per un itinerario di viaggio deve essere ottimizzato rispetto al tempo di percorrenza.
\end{enumerate}

\section{Criteri di successo}
\begin{itemize}
    \item Aumento del fatturato di almeno il 10\% entro un anno dal rilascio dell'applicazione, rispetto al fatturato
    dell'anno 2025;
    \item ottenere una valutazione media di almeno 7 punti su 10 da parte degli utenti entro un anno dal rilascio
    dell'applicazione, considerando almeno 100 voti. La valutazione viene richiesta a tutti gli utenti almeno una volta
    entro un anno dal rilascio dell'applicazione, tramite indagine interna all'applicazione, e se ne calcolerà la
    media al termine della raccolta.
\end{itemize}

\section{Scelte}

\subsection{Scelte architetturali}
\begin{itemize}
    \item In riferimento alla CoS 1, sviluppare prima una web app, e successivamente un'app mobile come wrapper
    dell'applicazione web;
    \item utilizzare un'architettura a microservizi per il backend, sfruttando Docker per la containerizzazione
    dei microservizi e Kubernetes per il deploy e l'orchestrazione dei container;
    \item ogni microservizio dovrà adottare un'architettura esagonale;
\end{itemize}

\subsection{Scelte tecnologiche}
\begin{itemize}
    \item dopo aver preso in considerazione database relazionali e documentali, si è deciso di utilizzare MongoDB come
    database documentale per la sua flessibilità e scalabilità, che si adattano meglio alle esigenze dell'applicazione,
    oltre ad essere il più conosciuto tra i membri del team;
    \item lo sviluppo deve seguire un approccio di Test-Driven Development (TDD);
    \item utilizzare Kotlin come linguaggio di programmazione per il backend, e Vue.js sfruttando TypeScript per il
    frontend della web app;
    \item in riferimento all'obiettivo 3, come metodo di associazione tra un'auto registrata e una colonnina di
    ricarica, si è deciso di utilizzare un codice QR, in quanto già presente nel software delle colonnine di ricarica;
    \item in riferimento all'obiettivo 1, l'autenticazione degli utenti avverrà tramite email e password, con
    l'adozione dell'autenticazione a due fattori tramite l'invio di un codice temporaneo via email.
\end{itemize}

\end{document}